\section{Engineering Roadmap}
\label{sec:roadmap}

\subsection{Structure}

The programme is organized into six phases across twenty-four months. Phases
are sequenced so that the highest-uncertainty work --- the imaging model ---
completes earliest, and so that a demonstrable artifact exists at the end of
every phase.

The governing principle is that \textbf{the imaging model is the risk and the
product is the schedule}. If the density model, print transfer, grain, and
halation are right, the remaining work is well-understood application
engineering with predictable duration. If they are wrong, no amount of interface
polish rescues the product. Phase 1 therefore front-loads the science, and its
exit criterion is a numerical one.

\subsection{Phase Definitions}

\paragraph{Phase 0 --- Foundations (months 1--2)}
Repository, module structure per Table~\ref{tab:layers}, CI, device farm,
\texttt{EmulsionCore} skeleton with the parameter model, and a minimal Metal
harness that renders a single pass to screen. Deliverable: an application that
decodes a DNG and displays it scene-referred with a working exposure slider.
Exit criterion: the null test (AC-7) passes.

\paragraph{Phase 1 --- The Imaging Model (months 2--7)}
The scientific core. Characteristic curve, development modulation, print
transfer, printer lights, interlayer inhibition, grain, halation. Reference CPU
implementation alongside every shader. Profile fitting tooling. Three negative
stocks and one print stock, fitted and validated.

Deliverable: a command-line renderer and a bare-bones application that produce
validated output. Exit criteria: AC-1 through AC-6 pass for the three stocks;
first perceptual evaluation round completed with cue list produced.

This is the phase where schedule risk concentrates, and it is deliberately given
six months for what an optimistic estimate would put at three.

\paragraph{Phase 2 --- Render Graph and Performance (months 6--10)}
Overlaps Phase 1 by one month. Graph compilation, fusion, aliasing, dirty
tracking, LUT baking, tile memory, progressive refinement, ROI rendering, export
tiling. Deliverable: NFR-01 through NFR-03 met on the target device set with the
Phase 1 model.

\paragraph{Phase 3 --- Application and Interface (months 8--15)}
Overlaps Phase 2 by two months. Capture subsystem, library, persistence,
edit history, the three workspaces, all panels, comparison modes, pipeline
inspector, export, accessibility. Deliverable: feature-complete against FR-01
through FR-22.

\paragraph{Phase 4 --- Stock Library and Refinement (months 13--19)}
Overlaps Phase 3. Remaining seven negative stocks and three print stocks fitted
and validated. Optical simulation, aging module. Second and third perceptual
evaluation rounds with iteration. Onboarding and the explanation content for
every control. Deliverable: content-complete.

\paragraph{Phase 5 --- Hardening and Launch (months 18--24)}
Overlaps Phase 4. Beta programme (target 500 external testers across the three
personas), crash and performance triage to NFR-10, localization, App Store
materials, pricing and subscription infrastructure, launch. Deliverable: 1.0 on
the App Store.

\begin{figure}[htbp]
\centering
\begin{tikzpicture}[font=\scriptsize, xscale=0.34, yscale=0.62]
% axis
\draw[emLine!40] (0,0) -- (24,0);
\foreach \m in {0,3,...,24} {
  \draw[emLine!40] (\m,0) -- (\m,-0.18);
  \node[below, font=\tiny] at (\m,-0.2) {M\m};
}
\foreach \m in {0,3,...,24} \draw[emLine!12] (\m,0) -- (\m,7.4);

% bars: name / start / end / row / color
\foreach \n/\s/\e/\r/\c in {
  {Ph.0 Foundations}/0/2/6.7/emLine,
  {Ph.1 Imaging model}/2/7/5.7/emAcc,
  {Ph.2 Graph + perf.}/6/10/4.7/emBlue,
  {Ph.3 App + UI}/8/15/3.7/emGreen,
  {Ph.4 Stocks + refine}/13/19/2.7/emKey,
  {Ph.5 Harden + launch}/18/24/1.7/emLine}
{
  \draw[fill=\c!25, draw=\c!80, rounded corners=1pt] (\s,\r) rectangle (\e,\r+0.6);
  \node[right, font=\tiny, text=black] at (\s+0.15,\r+0.3) {\n};
}

% milestones
\foreach \m/\lbl/\y in {7/{AC-1..6 pass}/5.7, 10/{NFR-01..03 met}/4.7,
                        15/{feature complete}/3.7, 19/{content complete}/2.7,
                        24/{1.0 ship}/1.7}
{
  \node[diamond, draw=emAcc, fill=emAcc!60, inner sep=1.1pt] at (\m,\y+0.3) {};
}
\node[font=\tiny, anchor=west] at (0.2,0.9) {$\diamond$ = phase exit milestone};
\end{tikzpicture}
\caption{Phase schedule. Overlaps are intentional: Phase 2 begins while the imaging model is still being fitted, because graph work depends on the pipeline's \emph{shape}, which stabilizes early, not on its parameters.}
\label{fig:gantt}
\end{figure}

\subsection{Milestones and Exit Criteria}

\begin{table}[htbp]
\caption{Milestones}
\label{tab:milestones}
\begin{center}
\renewcommand{\arraystretch}{1.15}
\begin{tabularx}{\columnwidth}{@{}ll>{\raggedright\arraybackslash}X@{}}
\toprule
\textbf{M} & \textbf{Milestone} & \textbf{Exit criterion} \\
\midrule
2  & Foundations & AC-7 null test passes; CI green; device farm operational. \\
7  & Model proven & AC-1--AC-6 pass on 3 stocks; perceptual round 1 complete; discrimination rate recorded. \\
10 & Interactive & NFR-01, 02, 03 met on A15 and A17 Pro with all stages enabled. \\
15 & Feature complete & All MUST requirements implemented; accessibility audit passed. \\
19 & Content complete & 10 negative + 4 print stocks validated; all explanation content written. \\
21 & Beta & 500 testers; crash-free $\geq 99.0\%$; perceptual round 3 discrimination $< 70\%$. \\
24 & Launch & NFR-10 met; App Store review passed; day-one support ready. \\
\bottomrule
\end{tabularx}
\end{center}
\end{table}

\subsection{Team Composition}

\begin{table}[htbp]
\caption{Staffing Plan (headcount by phase)}
\label{tab:staffing}
\begin{center}
\renewcommand{\arraystretch}{1.15}
\begin{tabular}{@{}lcccccc@{}}
\toprule
\textbf{Role} & \textbf{P0} & \textbf{P1} & \textbf{P2} & \textbf{P3} & \textbf{P4} & \textbf{P5} \\
\midrule
Imaging scientist        & 1 & 2 & 1 & 1 & 2 & 1 \\
GPU / graphics engineer  & 1 & 1 & 2 & 1 & 1 & 1 \\
iOS application engineer & 1 & 1 & 1 & 3 & 3 & 3 \\
Product designer         & --- & 1 & 1 & 1 & 1 & 1 \\
QA / test engineer       & --- & --- & 1 & 1 & 1 & 2 \\
Colorist consultant (PT) & --- & 0.5 & 0.5 & 0.5 & 1 & 0.5 \\
Product lead             & 1 & 1 & 1 & 1 & 1 & 1 \\
\midrule
\textbf{Total FTE} & \textbf{4} & \textbf{6.5} & \textbf{7.5} & \textbf{8.5} & \textbf{10} & \textbf{9.5} \\
\bottomrule
\end{tabular}
\end{center}
\end{table}

The imaging scientist role is the hard hire and the schedule's critical
dependency. The required profile --- comfort with sensitometry and color science
\emph{and} the ability to implement and validate numerically --- is rare, and
the recruiting process should begin before Phase 0. The colorist consultant is
part-time throughout and is the product's connection to professional practice;
this role is what keeps the vocabulary honest.

\subsection{Cost Model}

\begin{table}[htbp]
\caption{Cost Estimate (USD, fully loaded)}
\label{tab:cost}
\begin{center}
\renewcommand{\arraystretch}{1.15}
\begin{tabular}{@{}lrr@{}}
\toprule
\textbf{Category} & \textbf{Year 1} & \textbf{Year 2} \\
\midrule
Engineering and design salaries & 1{,}420{,}000 & 1{,}760{,}000 \\
Colorist and expert consulting & 60{,}000 & 75{,}000 \\
Reference material (film, processing, scanning) & 45{,}000 & 30{,}000 \\
Densitometry and measurement equipment & 38{,}000 & 6{,}000 \\
Device farm and CI infrastructure & 34{,}000 & 22{,}000 \\
Software, tooling, developer programme & 18{,}000 & 18{,}000 \\
Legal (trademark review, privacy) & 35{,}000 & 25{,}000 \\
Beta programme and user research & 15{,}000 & 55{,}000 \\
Marketing and launch & --- & 240{,}000 \\
Contingency (12\%) & 200{,}000 & 267{,}000 \\
\midrule
\textbf{Total} & \textbf{1{,}865{,}000} & \textbf{2{,}498{,}000} \\
\bottomrule
\end{tabular}
\end{center}
\end{table}

The reference material line funds the acquisition, exposure, processing, and
drum-scanning of physical film across the shipping stock list, plus the
densitometry to measure it. It is small in absolute terms and it is the line item
that most directly determines whether AC-1 through AC-4 can be met against real
data rather than against digitized datasheet curves. It should not be cut.

\subsection{Critical Path}

The critical path runs: imaging scientist hire $\to$ characteristic curve and
print transfer $\to$ profile fitting tooling $\to$ stock library $\to$
perceptual validation $\to$ launch. Every item on it is in the imaging domain.

Three float-heavy activities can absorb schedule pressure without moving launch:
optical simulation (FR-10, SHOULD), aging module (FR-11, SHOULD), and the
pipeline inspector (FR-20, SHOULD). These are the designated cut candidates and
are marked as such in the requirements table, so that a Phase 4 slip has a
prepared response rather than an improvised one.

\subsection{Post-Launch Sequence}

\begin{table}[htbp]
\caption{Post-1.0 Release Sequence}
\label{tab:postlaunch}
\begin{center}
\renewcommand{\arraystretch}{1.15}
\begin{tabularx}{\columnwidth}{@{}ll>{\raggedright\arraybackslash}X@{}}
\toprule
\textbf{Rel.} & \textbf{Target} & \textbf{Content} \\
\midrule
1.1 & +3 mo & Stability, stock additions, user-requested workflow fixes. \\
1.2 & +6 mo & Batch processing, tethered export presets, expanded sidecar interop (CDL, CUBE export of the baked pointwise chain). \\
2.0 & +12 mo & macOS companion sharing \texttt{EmulsionCore}; iCloud recipe sync. \\
2.5 & +18 mo & Video capture with real-time emulation; Apple Log input; gate weave and breathing. \\
3.0 & +24 mo & Node-based pipeline editing; user-authored stock profiles; profile marketplace. \\
\bottomrule
\end{tabularx}
\end{center}
\end{table}

The ordering is deliberate. macOS precedes video because it reuses existing code
and reaches the P2 persona, whereas video is a substantial new engineering
programme (temporal grain coherence, real-time constraint, codec integration).
Node-based editing comes last because it is the feature that most benefits from
an established user base to design against, and because the graph infrastructure
it requires is already built (Section~\ref{sec:rendergraph}).
